%-*-latex-*-
%-*-latex-*-
\newcommand\COURSE{ciss350}
\newcommand\ASSESSMENT{q10702}
\newcommand\ASSESSMENTTYPE{Quiz}
\newcommand\POINTS{\textwhite{xxx/xxx}}

\input{myquizpreamble}
\input{yliow}
\input{\COURSE}
%-*-latex-*-

%-*-latex-*-

%-*-latex-*-
\renewcommand\TITLE{\ASSESSMENTTYPE \ \ASSESSMENT}

\renewcommand\EMAIL{}
\renewcommand\AUTHOR{eagraham1@cougars.ccis.edu}

\textwidth=6in
\begin{document}
\topmatterthree

You are given the following binary tree node class:
\begin{Verbatim}[frame=single,fontsize=\footnotesize]
class BTNode
{
public:
    int key_;
    BTNode * parent_, * left_, * right_;
};
\end{Verbatim}
The following questions uses this class.
Make sure you test your code.

%------------------------------------------------------------------------------
\nextq
In the following tree
\begin{console}[fontsize=\footnotesize]
        4
       / \
      3   2
         / \
        0   5
       / \   \
      1   7   6
         /
        9
\end{console}
assume that the tree contains unique values.
Suppose \verb!p! points to the node with value 4.
Then the tree that \verb!p! points to (i.e., the root of the tree
that \verb!p! points to) has height 4.
If \verb!p! points to the node with value 2.
Then the tree that \verb!p! points to has height 3.
Complete the following function that returns the height
as described above.
\\
\ANSWER
\begin{answercode}
int height(BTNode * p)
{
  if(p == NULL)
      return 0;

  int leftH = height(p->left_);
  int rightH = height(p->right_);

  return 1 + (leftH > rightH ? leftH : rightH);
}
\end{answercode}

%------------------------------------------------------------------------------
\nextq
In the following tree
\begin{console}[fontsize=\footnotesize]
        4
       / \
      3   2
         / \
        0   5
       / \   \
      1   7   6
         /
        9
\end{console}
assume that the tree contains unique values.
\begin{enumerate}[nosep]
\li
Suppose \verb!p! points to the node with value 4
and \verb!q! points to the node with value 5.
Then we say that 5 is at a depth of 2 below 4
and when I call \verb!depth(p, q)!, I get 2.
\li
If \verb!p! points to the node with value 2
and \verb!q! points to the node with value 9.
Then we say that 9 is at a depth of 3 below 2
and when I call \verb!depth(p, q)!, I get 3.
\li
If \verb!p! points to the node with value 2
then \verb!depth(p, p)! returns 0.
\end{enumerate}
The above describes the cases when
\verb!p! and \verb!q! points to nodes
and either \verb!p! and \verb!q! are the same
or \verb!*q! is a descendent of \verb!*p!.
In all other cases, the function should return $-1$.
Complete the following function that returns the depth
as described above.
\\
\ANSWER
\begin{answercode}
int depth(BTNode * p, BTNode * q)
{
  if(p == NULL && q == NULL
      return -1;
  if(p == q)
      return 0;

  BTNode * curr = q;
  int d = 0;

  while(curr != NULL)
  {
      if(curr == p)
         return d;

      curr = curr->parent_;
      ++d;
  }

  return -1;
}
\end{answercode}

%------------------------------------------------------------------------------
\nextq
In the following tree
\begin{console}[fontsize=\footnotesize]
        4
       / \
      3   2
         / \
        0   5
       / \   \
      1   7   6
\end{console}
assume that the tree contains unique values.
\begin{enumerate}[nosep]
\li
Suppose \verb!p! points to the node with value 4,
the depth of 0 (below 4) is said to be 2.
The depth of 6 (below 4) is 3.
\end{enumerate}
Check the notes.
Complete the following function that returns the depth
of the node with \verb!key_! value of \verb!target!
below the node that \verb!p! points to.
If \verb!target! cannot be found, \verb!-1! is returned.
\\
\ANSWER
\begin{answercode}
int depth(BTNode * p, int target)
{
  if(p == NULL)
    return -1;

   std::queue<BTNode*> nodeq;
   std::queue<int> depthq;

   nodeq.push(p);
   depthq.push(0);

   while(!nodeQ.empty())
   {
     BTNode * curr = nodeq.front();
     nodeq.pop();

     int d = depthq.front();
     depthq.pop();

     if (curr->key_ == target)
       return d;

     if (curr->left_ != NULL)
     {
       nodeQ.push(curr->left_);
       depthQ.push(d + 1);
     }

     if (curr->right_ != NULL)
     {
       nodeQ.push(curr->right_);
       depthQ.push(d + 1);
     }
   }

   return -1;
}
\end{answercode}
(Hint: Do a BF traversal. Instead of putting node
pointers into the queue, you want to put the node pointer together
with its depth into the queue.
Note that even if the tree has repeat \verb!key_! values,
if you do a BF traversal,
the depth computed is the \textit{shallowest} depth of all the nodes with the
\verb!key_! value of \verb!target!.)

%------------------------------------------------------------------------------
\nextq
Complete the function below that returns true if
\verb!p! and \verb!q! point to roots of binary trees which are the same.
In this case, same means the tree has the same structure, i.e.,
the root nodes have the same \verb!key_! value,
while left child nodes have the same value, then right child node have the
same \verb!key_! value, etc.
And if a node in the tree \verb!p! points to
does not have a left child, then the corresponding node in the tree \verb!q!
points to does not have a left child; likewise for right child.
The following trees are the same
\begin{console}[fontsize=\footnotesize]
        4             4
       / \           / \
      3   2         3   2
         / \           / \
        0   5         0   5
\end{console}
The following trees are different:
\begin{console}[fontsize=\footnotesize]
        4             4
       / \           / \
      3   2         3   2
         / \           / \
        0   5         5   0
\end{console}
The following trees are different:
\begin{console}[fontsize=\footnotesize]
        4             4
       / \           / \
      3   2         3   2
         / \             \
        0   5             5
\end{console}
Don't forget to check the parent pointer.
\\
\ANSWER
\begin{answercode}
bool is_same_tree(BTNode * p, BTNode * q)
{
  if(p == NULL && q == NULL) return true;
  if(p == NULL || q == NULL) return false;
  if(p->key_ != q->key_) return false;

  if ((p->parent_ == NULL && q->parent_ != NULL) ||
     (p->parent_ != NULL && q->parent_ == NULL))
    return false;

  if (p->parent_ != NULL && q->parent_ != NULL &&
      p->parent_->key_ != q->parent_->key_)
    return false;

  return is_same_tree(p->left_, q->left_) &&
         is_same_tree(p->right_, q->right_);
}
\end{answercode}

%------------------------------------------------------------------------------
%-*-latex-*-
\newpage
\input{instructions}
\end{document}

